\section{蒸馏：把稀疏 teacher 转换为更易部署的 student}

MoE 的训练优势不等于所有部署都适合保留专家分布。低 batch、边缘设备或参数存储受限场景可能更需要一个较小 dense student。Distillation（知识蒸馏）让 student 学习 teacher 的软分布、隐藏表示或任务输出，但它会把“巨大条件参数池”压缩回更小容量，因此必然是质量、参数、延迟与实现复杂度之间的转换。

\lecturefigure{slide-33-distillation-techniques.jpg}{Switch Transformer 的蒸馏技术与 DistilBERT 风格基线。}{Stanford Online 官方视频 00:47:42--00:48:23。}

一个常见蒸馏目标可写为：

\[
\mathcal{L}=\lambda\mathcal{L}_{\text{hard}}+(1-\lambda)\tau^2
\operatorname{KL}\!\left(p_T^{(\tau)}\,\|\,p_S^{(\tau)}\right).
\]

其中，$\mathcal{L}_{\text{hard}}$ 是 student 对真实标签或原始预训练目标的 loss；$p_T^{(\tau)}$ 与 $p_S^{(\tau)}$ 分别是 teacher 和 student 在温度 $\tau$ 下的软分布；KL 是分布差异；$\lambda$ 平衡硬目标与 teacher 信号。不同蒸馏方案还可加入 hidden-state matching 或仅使用 teacher 生成数据。

\begin{knowledgebox}{术语消化：蒸馏不是量化}
Distillation 训练一个新 student 去逼近 teacher 行为；quantization 则降低同一模型权重或激活的数值位宽。前者改变参数结构和容量，后者主要改变表示精度。两者可以组合，但失败模式不同。
\end{knowledgebox}

\subsection{蒸馏预训练模型：压缩参数，也压缩条件容量}

预训练蒸馏比较 dense、sparse teacher 与不同 student。读表时应同时看参数、训练或推理 FLOPs、上游质量和恢复比例。若 student 参数大幅下降但只保留部分 teacher 质量，这不是蒸馏失败，而是明确的 Pareto 交易；真正问题是它是否比直接训练同规模 dense model 更好。

\lecturefigure{slide-34-distilling-pretrained-model.jpg}{预训练阶段的 sparse teacher 蒸馏与参数/质量权衡。}{Stanford Online 官方视频 00:48:24--00:50:58。}

\begin{knowledgebox}{读表：必须加入“直接训练同规模 student”基线}
只比较 teacher 与 distilled student，会把“容量变小后的必然下降”误当作蒸馏结论。应加入从头训练的同结构 student，判断 teacher 信号究竟恢复了多少质量，以及额外蒸馏成本是否值得。
\end{knowledgebox}

\teachervoice{Q\&A 中的实质问题是部署目标：若最终必须得到小 dense model，为什么不一开始就训练它？蒸馏只有在 sparse teacher 能提供更好的软目标、数据扩增或表示监督，并且这些收益超过 teacher 训练成本时才成立。}

\subsection{蒸馏 fine-tuned SuperGLUE 模型}

另一条路线是先让 Switch teacher 在 SuperGLUE 上完成 fine-tuning，再蒸馏任务行为。它更贴近特定产品，但泛化范围更窄：student 可能保留任务分数，却不再拥有 teacher 的通用预训练能力。部署前应决定需要“通用 student”还是“任务专用 student”。

\lecturefigure{slide-35-distilling-finetuned-superglue.jpg}{蒸馏已 fine-tune 的 SuperGLUE teacher：任务质量与模型规模的转换。}{Stanford Online 官方视频 00:50:59--00:57:07。}

\begin{importantbox}{蒸馏验收的四个维度}
同时报告：相对 teacher 的质量保留、相对同规模 scratch baseline 的增益、端到端训练成本、以及目标硬件上的延迟/吞吐。若只报告参数压缩比，就无法判断部署是否真的更优。
\end{importantbox}

\begin{warningbox}{蒸馏不能消除 teacher 的偏差}
Student 会学习 teacher 的分布与错误。若 sparse router 对语言、主题或群体存在不均衡，蒸馏可能把这种行为固化进 dense student，而不是自动平均掉。
\end{warningbox}

\subsection{本章小结}

蒸馏是从高容量 sparse teacher 到小型部署模型的有损转换。正确基线是同规模直接训练 student，正确指标同时包含质量、训练成本和真实硬件性能。

